\begin{claim}{F.20}{Proved}{An exact cross-divisor action and its finite wreath images}
Retain the labelled source, the two-track group, and the subgroups
$K^+\leq D^+$ of F.17--F.19. Select the cross divisor
between the old point $v_1$ and the positive new point $a$.
Its full relative path coordinate is
$\mathcal F=\widehat F(x,y,z)$, with the three specified
meridians corresponding, in order, to $0,v_\zeta,1$.
The actual source $L=\widehat{\pi_1(Y)}=\widehat F_{11}$
has the continuous action
\[
       A_g(v)=q_g^{-1}\alpha_g(v)p_g
                 \qquad(g\in L,\ v\in\mathcal F),
                                                               \tag{F.20.1}
\]
where $\alpha_g,p_g,q_g$ are the literal based braid data
specified below. The diagonal-coset map proving this
formula, its inverse, and every based source transport
are retained in the proof.

Let $\mathscr Q$ be the following explicitly enumerated
finite quotients of this same $\mathcal F$:
the $82$ source-invariant $A_5$ quotients in the enumeration
below, all $3,321$ products of two distinct such quotients,
and the six mod-two homology-cover quotients
\[
  Q_\rho=\mathcal F/
       \overline{[N_\rho,N_\rho]\,N_\rho^2},
       \qquad N_\rho=\ker(\rho:\mathcal F\to A_5),
       \qquad |Q_\rho|=60\cdot2^{121},
                                                               \tag{F.20.2}
\]
whose four peripheral images under $\rho$ are all nonidentity.
Every action (F.20.1) on every member of $\mathscr Q$ is
transitive. In each of the six last cases the image of the
source filling kernel acts by right translations through
a subgroup of index four; the remaining actual source
generators connect all four cosets.

For each $Q\in\mathscr Q$ and each prime $\ell$ there is
an explicitly specified continuous homomorphism
\[
 \chi_{Q,\ell}:D^+\longrightarrow\mathbb F_\ell[Q]
 \quad\hbox{with}\quad
 \chi_{Q,\ell}(K^+)=\chi_{Q,\ell}(D^+)=\mathbb F_\ell[Q].
                                                               \tag{F.20.3}
\]
It is obtained from the actual one-point forgetting map
and the finite wreath assignment in (F.20.15) below.
Thus these maps give the same full finite images for
the two specified subgroups. This conclusion does not
identify $K^+$ with $D^+$ or solve the remaining
source-membership equation of F.19.
\end{claim}

\begin{proof}
\emph{The selected two-track space and its actual diagonal.}
Normalize the common labels $0,\infty,1$ as in F.19 and put
\[
 \begin{gathered}
 \zeta=e^{2\pi i/5},\qquad
 Y=\mathbb P^1\setminus\{0,\infty,\pm\zeta^j:0\leq j<5\},\\
 a(s)=\frac{\zeta-s}{\zeta+s},\quad
 B(s)=a(s)^2,\quad
 P(s)=\left(\frac{1-s}{1+s}\right)^2,\quad Q(s)=a(s).
 \end{gathered}                                                \tag{F.20.4}
\]
The letter $Q(s)$ in this display is a point coordinate;
a finite quotient $Q\in\mathscr Q$ will always be named as such.
Forget the other three old labels and the negative new label
in the F.19 two-track space $X$. The target is exactly
\[
 X_1=\{(B,U,V):B\ne0,1,\quad
                  U,V\notin\{0,1,B\}\}
       =\mathcal M_{0,5}\times_{\mathcal M_{0,4}}\mathcal M_{0,5}.
                                                               \tag{F.20.5}
\]
It allows $U=V$; that equality is its selected diagonal.
The curve map is literally $s\mapsto(B(s),P(s),Q(s))$.
The selected divisor upstairs in $X$ is its inverse image.

The comparison with the original F.17 basepoint is
also specified. Write
$s_0=i(\theta^2-2)\zeta$, as in F.12, and take the
polygonal path with successive vertices
\[
        s_0,\quad 3s_0,\quad 3s_0+6i,\quad
                         6+6i,\quad 2+2i.                    \tag{F.20.4a}
\]
It lies in $Y^\dagger$.
Indeed $|s_0|^2=\theta^2-1>1$.
For the four $A_u$ of F.17.2, the imaginary parts
of $A_u\overline{s_0^2}$ are respectively
$\theta^3/2,-\theta^3/2,
 (5\theta-3\theta^3)/2,
 (-5\theta+3\theta^3)/2$ in the order
$u=1,\zeta^2,\zeta^3,\zeta^4$; all are nonzero.
Thus its first radial segment misses the outer
collision roots. The inner collision roots have
absolute value less than one.
All later segments lie outside the disk of radius
two: $3\operatorname{Re}s_0<-2$, the middle
horizontal-direction segment has imaginary part
at least six, and the last segment has both
coordinates at least two.
Every collision root has absolute value less
than two, by $|A_u|<3$ and $|B_u|<1$.
The exact producer checks these inequalities.

Call this path $\eta$.
The based maps are $g\mapsto\eta g\eta^{-1}$
on the source and conjugation by the actual
labelled configuration path $\phi^\dagger(\eta)$
on the ten-point target, with the identical
forgettings of that path on every smaller target.
The composites commute because these are
literal paths of the displayed rational labels.
They carry the $K^+,D^+$ at $s_*=2+2i$ to the
original F.19 subgroups at $s_0$.
The homomorphisms in (F.20.3) on those original
subgroups include the inverse of this specified
basepoint change.

The forgetting projection is a configuration fibration with
connected fibers. Over its diagonal the additional three
old points and the other new point still vary in their
respective connected configuration fibers. Consequently
the image of the divisor fundamental group is the full
inverse image of the diagonal fundamental group: lift a
diagonal loop through this fibration, and correct its
endpoint within the connected fiber. The kernel of the
projection is represented by loops in that same fiber
and already lies in the divisor image.

The same comparison holds on full completions.
For clarity, write $P_j=\widehat{\pi_1(\mathcal M_{0,j})}$
with the specified label maps. The diagonal inverse image
is the fiber product $P_8\times_{P_5}P_6$; its kernel over
$P_8$ is the actual one-point-pushing $\widehat F_4$.
The discrete sequence has an actual continuous
configuration section, obtained by placing the forgotten
point on the positive real axis beyond
$2+|B|+|U|$, and then choosing the based fiber path.
The characteristic finite kernel quotients used in F.19
give its full induced topology and the explicit completed
semidirect-product inverse maps. They identify its
completion with this stated fiber product.
The identical argument for (F.20.5) gives
$\mathcal Q_1=P_5\times_{P_4}P_5$.
Thus the full coset space for the selected divisor in
F.19 is carried homeomorphically onto
$\mathcal Q_1/\operatorname{diag}(P_5)$.
No closure/intersection exchange is used.

\emph{The based source and the two reference transports.}
Use source basepoint $s_*=2+2i$, reference point $R=8+3i$,
and projection $z\mapsto\operatorname{Re}((1+371i/1000)z)$.
For each of the eleven centers
\[
       0,\ 1,\zeta,\zeta^2,\zeta^3,\zeta^4,\
          -1,-\zeta,-\zeta^2,-\zeta^3,-\zeta^4
                                                               \tag{F.20.6}
\]
take the straight whisker from $s_*$ to the northeast
corner of the axis-parallel square of half-side
$31/1000$ about that center, traverse the square
counterclockwise, and return along that same whisker.
Call these loops $g_0,\ldots,g_{10}$ in the displayed order.
The certificate verifies that the squares are pairwise
disjoint, each whisker misses all other squares, and
the whiskers have distinct directions from their common
basepoint. The basepoint lies northeast of every square,
so its own whisker first meets that square at its corner.
These squares and the embedded star give an actual
geometric free basis of $\pi_1(Y)$.

At $s_*$, transport the fourth finite point along the
straight segment from $R$ to $P(s_*)$ for the old track
and to $Q(s_*)$ for the new track, keeping $0,1,B(s_*)$
fixed. Both paths avoid those three punctures; their
exact cross-product checks are retained. Call these
paths $T_o,T_n$. A source loop gives the two based
four-strand braids
$T_o\,P(g)\,T_o^{-1}$ and $T_n\,Q(g)\,T_n^{-1}$.
The notation here retains the actual intermediate
three-strand path $B(g)$.
At the common reference configuration the projected
order is exactly $(0,B(s_*),1,R)$.
Use the same disk basepoint and straighten the four
points within their projection-order chamber. A generator
$\sigma_j$ is the literal clockwise exchange of adjacent
points: the left strand passes above the right strand.
For example, after translation and positive scaling, the
two paths
\[
 2t-1+4it(1-t),\qquad -2t+1-4it(1-t)
       \quad(0\leq t\leq1)
\]
have exactly this exchange and the producer's positive
crossing. Their squared distance from the midpoint is
$((2t-1)^2-1/2)^2+3/4$, so the pair never collides.
These geometric generators give the Artin presentation
of the actual braid group. The symbols $x,y,z,d$ below
are its ordered free coordinates; the exact identification
of $x,y,z$ with the actual point-pushing fibre, including
its inverse based loops, is given by (F.20.8)--(F.20.9).

\emph{The Artin coordinate and the diagonal-coset formula.}
Use the faithful Artin substitutions
\[
 \sigma_j:
 X_j\longmapsto X_jX_{j+1}X_j^{-1},\quad
 X_{j+1}\longmapsto X_j ,
 \qquad \sigma_j(X_k)=X_k\ (k\ne j,j+1),
                                                               \tag{F.20.7}
\]
in the word-composition convention of \cite{A-Artin1947}.
Path products are read from left to right in time, and
the substitution for a displayed word $b_1\cdots b_r$
is $\phi_{b_1}\circ\cdots\circ\phi_{b_r}$, as in the
producer. The displayed substitutions and their inverses
satisfy the Artin relations, so the named geometric
generators define this homomorphism into the free-group
automorphisms. The following fibre-coordinate comparison
determines its actual point-pushing action.

For a pure four-strand braid $\beta$, write
$\phi_\beta(d)=W_\beta dW_\beta^{-1}$ and let $w_\beta$
be the word obtained by filling $d$ in $W_\beta$.
The ambiguity in $W_\beta$ is a power of $d$ and
therefore disappears under this filling.
Write $\alpha_\beta$ for the filled automorphism of
$F(x,y,z)$. Direct substitution gives
\[
 w_{\beta\gamma}=\alpha_\beta(w_\gamma)w_\beta,\qquad
 u_\beta=w_\beta^{-1},\qquad
 u_{\beta\gamma}=u_\beta\alpha_\beta(u_\gamma).
                                                               \tag{F.20.8}
\]
To identify the actual point-pushing kernel, let
$\beta_j=\sigma_3\cdots\sigma_{j+1}\sigma_j^2
 \sigma_{j+1}^{-1}\cdots\sigma_3^{-1}$.
The first descending word carries the last strand along
the lower-half-plane whisker, the square makes a clockwise
lasso about puncture $j$, and the inverse word returns
that whisker. Deleting the last strand contracts the
old motion within its projection-order chambers. Thus
these are the three actual lower-whiskered clockwise
lassos, a free basis of the punctured-plane fibre.
Direct substitution gives their $u$-images
\[
        z^{-1}y^{-1}x^{-1}yz,\qquad
        z^{-1}y^{-1}z,\qquad z^{-1}.
                                                               \tag{F.20.9}
\]
This triangular substitution is its own inverse, by
free reduction. The inverse map sends $x,y,z$ to the
same three words with $\beta_1,\beta_2,\beta_3$
substituted for $x,y,z$. These are specified conjugates
of the inverse clockwise lassos, so they are actual
based puncture meridians. Both composites are the
identity on the named free bases. This gives the actual
fibre isomorphism and then its full profinite inverse.
Adding the last strand gives the standard section on
the three-strand braid group, with $u=1$.
The cocycle (F.20.8) therefore makes conjugation by that
section exactly the filled Artin substitution, and makes
conjugation by a general pure braid $b$ the map
$v\mapsto u_b\alpha_b(v)u_b^{-1}$.
Its restriction gives the section of $P_5\to P_4$
used in this coordinate. Filling a puncture kills its
lower-whiskered lasso; the same triangular formula
commutes with filling the corresponding variable and
retaining the order of the others. These are the actual
based filling maps in the displayed coordinates.

Here the normalized moduli groups have not been replaced
by braid groups without a map. The coordinate isomorphism
\[
 (z_0,z_1,z_B,z_R)\longmapsto
 \left(z_0,z_1-z_0,\
       \frac{z_B-z_0}{z_1-z_0},
       \frac{z_R-z_0}{z_1-z_0}\right)
                                                               \tag{F.20.10}
\]
identifies the ordered configuration space with
$\mathbb C\times\mathbb C^\times\times\mathcal M_{0,5}$;
its inverse multiplies the last two coordinates by
$z_1-z_0$ and adds $z_0$.
The analogous three-point formula has
$\mathcal M_{0,4}$ as last factor.
The original paths keep $z_0=0,z_1=1$.
The braid section commutes with forgetting to these
two labelled points, so it restricts to the kernels
of the displayed affine-coordinate projections.
The section and its retraction extend to full
completions and have the same composites.

After the based transports, both factors of
$\mathcal Q_1$ are this same $P_5$. The mutually inverse
coset maps are
\[
 \begin{aligned}
 \mathcal Q_1/\operatorname{diag}(P_5)&\longrightarrow\mathcal F,
       &(b_o,b_n)\operatorname{diag}(P_5)&\longmapsto b_nb_o^{-1},\\
 \mathcal F&\longrightarrow
                 \mathcal Q_1/\operatorname{diag}(P_5),
       &v&\longmapsto(1,v)\operatorname{diag}(P_5).
 \end{aligned}                                                \tag{F.20.11}
\]
The first value belongs to the kernel because both
factors have the same base image. Both maps are
continuous and their composites are the identity.
Left multiplication by the two source braids acts
by $v\mapsto b_nvb_o^{-1}$.
Substituting $b_o=(p_g^{-1},\alpha_g)$ and
$b_n=(q_g^{-1},\alpha_g)$ from (F.20.8) proves
(F.20.1), including its order of multiplication.

\emph{The complete generator data and exact crossing certificate.}
Set $P=xyz$, $k=Pxyx$, and $t=Pxyx^{-1}$.
Let $\tau_1,\tau_2$ be the automorphisms (F.20.7)
of $F(x,y,z)$ associated with $\sigma_1^2,\sigma_2^2$.
The complete source data are
\[
\begin{array}{c|c|c|c}
i&\alpha_{g_i}&p_{g_i}&q_{g_i}\\ \hline
0&\tau_1\tau_2^{-1}\tau_1^{-1}
  &(xy)x(xy)^{-1}P^{-1}&z^{-1}xy^{-1}x^{-1}\\
1&1&x^{-2}&1\\
2&\tau_1^{-2}&1&(xy)^{-1}\\
3&1&1&1\\
4&1&ky^{-1}k^{-1}&1\\
5&1&1&1\\
6&1&P^2&1\\
7&(\tau_1\tau_2)^2&tz^{-1}x^{-1}&t\\
8&1&1&1\\
9&1&xy^{-1}x^{-1}&1\\
10&1&1&1 .
\end{array}                                                   \tag{F.20.12}
\]
These are values in the specified basis, not
independently chosen conjugacy classes.

The executed producer
\path{code/part_a/cross_divisor_monodromy_f20.py}
works in $\mathbb Q(\theta)$, where
$\theta^4-5\theta^2+5=0$ and $19/10<\theta<2$, using
$\zeta=(\theta^2-3+i\theta)/2$.
Every sign is separated by rational interval arithmetic
at that particular real root.
For each rational configuration segment and each pair
of strands it multiplies the projected difference
numerator by the conjugate of its denominator.
Its real polynomial has exactly the projected
crossings; its imaginary polynomial gives their
heights, since the omitted denominator is positive.
All possible poles and collisions of either separate
track are among (F.20.6): in particular
$P(s)=B(s)$ has numerator
$4s(1-\zeta)(\zeta-s^2)$, whose remaining roots
are $\pm\zeta^3$.
The source-basis clearance therefore also checks that
these paths stay in the stated configuration spaces.

Exact Bernstein subdivision supplies the root count.
On a rational interval, multiplication by the positive
binomial coefficients changes its Bernstein sequence
to the coefficients of the usual fractional-linear
transform for Descartes' rule. Zero sign variation
excludes roots and one variation gives exactly one
root with multiplicity one. Intervals with larger
variation are subdivided. A subdivision point that
is an exact root is avoided by choosing another
rational fraction; a nonzero polynomial of degree
$d$ cannot vanish at $d+1$ such distinct choices.
The output retains the complete projected and height
polynomials, the terminal root-isolating intervals
and their Bernstein signs, and the counts of discarded
root-free leaves; the executed code reproduces the
subdivisions. The certificate further
separates a nonzero derivative and nonzero height
on every crossing interval.
Intervals for different pairs are disjoint.
Their exact order, initial projected order,
adjacency at each crossing, and final projected
order give the literal braid word.

There are $112$ certified forward segments,
$672$ projected-pair checks, $172$ crossings, and
$72,380$ exact number-field sign calls.
The return whisker is the literal inverse path,
so its braid is the inverse of its certified
forward braid. The two reference transports and
all eleven old and new based braids are included.
Deleting the fourth strand gives respectively the
three nontrivial base braid words
$\sigma_1^2\sigma_2^{-2}\sigma_1^{-2}$,
$\sigma_1^{-4}$, and
$(\sigma_1^2\sigma_2^2)^2$ at indices $0,2,7$.
The second producer verifies this strand deletion,
all entries of (F.20.12), and the inverse substitution
(F.20.9) by exact free reduction.
Neither producer reads a numerical discovery file.

\emph{The actual full one-point quotient and its finite wreath maps.}
Forget the negative new point in F.19's two-point
configuration kernel. The two Fadell--Neuwirth rows
and their filling maps are
\[
 \begin{array}{ccccccccc}
 1&\longrightarrow&\widehat F_8&\longrightarrow&J
                 &\longrightarrow&\widehat F_7&\longrightarrow&1\\
 &&\downarrow&&\downarrow&&\downarrow\\
 1&\longrightarrow&\widehat F_4&\longrightarrow&J_4
                 &\longrightarrow&\mathcal F&\longrightarrow&1 .
 \end{array}                                                  \tag{F.20.13}
\]
Here $J$ and $J_4$ are the actual eight- and four-puncture
two-point groups of F.19, and $D^+=\ker(J\to J_4)$.
Every vertical filling is onto. The full topologies
and exact rows are those proved in F.19 by the
actual configuration sections and finitely generated
free kernels.
Put $R_0=\ker(\widehat F_7\to\mathcal F)$.
The induced map
\[
                         D^+\longrightarrow R_0
                         \quad\hbox{is onto}.               \tag{F.20.14}
\]
Indeed lift an element of $R_0$ to $J$.
Its image in $J_4$ lies in the displayed $\widehat F_4$;
an element of $\widehat F_8$ has its inverse as image.
Multiplication by that element gives the required
lift in $D^+$ without changing its $\widehat F_7$ image.

Choose the geometric free basis of the seven-puncture
one-point group consisting of $x,y,z$ and the four old
puncture meridians $u_j$. Its filling kills the $u_j$
and retains $x,y,z$. The whisker of the selected
$u_1$ is chosen to project to the old reference
transport specified above after changing the
one-point basepoint by $T_n$; at the original
one-point basepoint its projected path is
$T_n^{-1}T_o$. The straight $T_n$ also avoids
all four other old points, as checked by their
nonzero exact cross products.
This is an explicit
allowable basis adjustment: a difference word $w$
in the retained $x,y,z$ replaces $u_1$ by
$wu_1w^{-1}$, fixes the other generators, and has
inverse $u_1\mapsto w^{-1}u_1w$.
The difference word is the class of the two stated
based paths. Thus the full filling and its splitting
are unchanged, with their exact maps retained.

For $\rho:\mathcal F\to Q$ in the stated collection
and a prime $\ell$, the finite free-group assignment is
\[
 \begin{aligned}
 \widehat F_7&\longrightarrow
                    \mathbb F_\ell[Q]\rtimes Q,\\
 x,y,z&\longmapsto(0,\rho(x)),(0,\rho(y)),(0,\rho(z)),\\
 u_1&\longmapsto(e_1,1),\qquad
 u_j\longmapsto(0,1)\quad(j\ne1).
 \end{aligned}                                                \tag{F.20.15}
\]
Here $Q$ acts by its actual left regular permutation
of the basis $e_a$, $a\in Q$.
On $R_0$ it is a homomorphism to the additive module
and is onto: conjugating $u_1$ by a retained word
with image $a$ gives $e_a$.
Composing with (F.20.14) defines $\chi_{Q,\ell}$
and proves its surjectivity on $D^+$.
This supplies continuity through an actual finite
quotient; it does not complete an infinitely generated
discrete meridian module without a topology argument.

The source filling $Y^\dagger\to Y$ kills the sixteen
specified positive collision meridians. Its image
$K^+$ is their closed normal closure inside the
drilled source image, as in F.17--F.19.
Forgetting the negative new point kills its eight
collision meridians. The assignment (F.20.15) kills
the six positive collisions with the other old labels.
The remaining two positive meridians have images
$e_{a_+},e_{a_-}$ for the two actual transported
relative path classes in $Q$.
Their exponent is one, by the simple holomorphic
collision and annular calculation of F.17.

Conjugation by a source lift transports these basis
elements by exactly (F.20.1).
One can see this either by (F.20.11) or directly:
a based old meridian is carried to its conjugate
by the old track, while the module basepoint is
carried by the new track. Their two transports
give $q_g^{-1}\alpha_g(v)p_g$ in that order.
Different lifts of the same element of $L$ differ
by an element of $R_0$, whose conjugation acts
trivially on this additive image.
Consequently, writing $\mathcal O_\pm$ for the
two actual source orbits of $a_\pm$, one has the
literal finite-image equality
\[
 \chi_{Q,\ell}(K^+)
       =\operatorname{span}_{\mathbb F_\ell}
                  \{e_a:a\in\mathcal O_+\cup\mathcal O_-\}.
                                                               \tag{F.20.16}
\]
The labels $a_\pm$ have not been moved to the
identity by an unproved normalization.
The transitivity proved next makes their particular
values unnecessary for (F.20.3).

\emph{The finite simple quotients and their paired actions.}
The executed producer
\path{code/part_a/cross_divisor_finite_f20.py}
reads only the preceding exact geometric output.
It constructs $A_5$ as the sixty even permutations
of five letters, with its exact multiplication
and inverse tables. Conjugation by $S_5$ gives
all its $120$ automorphisms: every automorphism
permutes the five Sylow four-groups, and this
action is faithful because the resulting faithful
$A_5$-action determines each inner automorphism.
The conjugacy-class sizes are checked as
$1,12,12,15,20$; no proper nontrivial union
including the identity has order dividing $60$.
This also verifies the simplicity input below.

The program enumerates all $200,160$ generating
triples and their $1,668$ automorphism orbits.
These are exactly the epimorphism kernels.
It evaluates each of the eleven specified
$\alpha_{g_i}$ on each triple and finds precisely
$82$ invariant kernels. For each one it records
the unique descending automorphism $\sigma_i$,
the two word values $\rho(p_{g_i}),\rho(q_{g_i})$,
and the sixty-point permutation
\[
       v\longmapsto
       \rho(q_{g_i})^{-1}\sigma_i(v)\rho(p_{g_i}).
                                                               \tag{F.20.17}
\]
Every one of these $82$ actions has a single orbit.

For two distinct kernels, the map to $A_5^2$
is onto. Indeed a proper subdirect subgroup of
two copies of a nonabelian simple group is the
graph of an automorphism: the two projection
kernels are normal, and simplicity leaves
only the full product or that graph.
The graph would identify the two kernels,
contrary to their enumeration.
The program applies each source generator
simultaneously to both coordinates; all
$3,321$ such actions have one orbit of size
$3,600$. This is a test of the actual joint
action, not two separate marginal counts.

\emph{The six noncentral homology-cover refinements.}
Exactly six of the invariant triples have
$\rho(x),\rho(y),\rho(z),\rho(xyz)$ all nonidentity.
For such a triple put $G=A_5$, and consider the
connected based Cayley covering graph of the
three-petalled rose. Its chain groups and
boundary over $\mathbb F_2$ are
\[
 C_1=\mathbb F_2[G]^3,\quad C_0=\mathbb F_2[G],
 \qquad \partial e_{i,g}=e_{g\rho(x_i)}-e_g .
                                                               \tag{F.20.18}
\]
For a word $w$, let $v(w)$ be its actual oriented
edge chain and $a(w)=\rho(w)$.
Concatenation gives
$v(ww')=v(w)+a(w)v(w')$.
Every chain with boundary $e_g-e_1$ is the chain
of a word ending at $g$: subtract one path to
$g$ and realize the remaining graph homology
class by based loops.
Thus (F.20.2) is exactly the finite group
\[
 \{(v,g)\in\mathbb F_2[G]^3\rtimes G:
                       \partial v=e_g-e_1\},
 \quad
 (v,g)(w,h)=(v+gw,gh).
                                                               \tag{F.20.19}
\]
Its kernel has dimension $180-59=121$.
This description does not assert that this
particular extension splits.

Each actual $\alpha_{g_i}$ preserves $N_\rho$
and its mod-two homology kernel.
The lifted action is computed by the edge
chains of its three actual generator images.
If $J_i$ is that semilinear substitution matrix,
$\sigma_i$ its permutation of $G$, and
$(v_p,a_p),(v_q,a_q)$ the two word values,
its affine action is literally
\[
 (v,a)\longmapsto
 \left(a_q^{-1}\bigl(J_i v-v_q+\sigma_i(a)v_p\bigr),
                         a_q^{-1}\sigma_i(a)a_p\right).
                                                               \tag{F.20.20}
\]
The code checks these formulas against direct
word evaluation in $660$ exact comparisons.

The filling $Y\to
\mathbb P^1\setminus\{0,\infty,\pm\zeta\}$
kills exactly $g_1,g_3,g_4,g_5,g_6,g_8,g_9,g_{10}$.
Their base automorphisms and new words in
(F.20.12) are the identity, so they act by
right translations through their old words.
Conjugating such a translation by a general
source action gives the right translation
\[
                         h\longmapsto
                  p_{g_i}^{-1}\alpha_{g_i}(h)p_{g_i}.
                                                               \tag{F.20.21}
\]
The finite subgroup $H_\rho$ generated by these
translations and their indicated conjugates is
computed exactly as follows. Choose representatives
of its projection into $G$. Schreier products
give vectors in $\ker\partial$; their span,
closed under the projection's regular action,
is precisely its kernel. Add the images of
the representatives and this kernel under
(F.20.21), and repeat until stable.
Each step is finite linear algebra and finite
group multiplication. Stability under the
given injective maps on a finite subgroup
also gives stability under their inverses.
This computes the full subgroup generated
by the named source filling kernel.

In all six cases its projection is all of
$A_5$ and its vector kernel has dimension
$119$. Hence $[Q_\rho:H_\rho]=4$.
The source action descends to its right
cosets, because
$A_g(vh)=A_g(v)\,
 (p_g^{-1}\alpha_g(h)p_g)$.
The program reduces those cosets by the
specified representatives and row-reduced
kernel basis, and finds one orbit of size
four in every case.
Right translations by $H_\rho$ are themselves
source actions, so this proves transitivity
on the entire group $Q_\rho$.

These computations prove transitivity for
every finite quotient in the statement.
Equation (F.20.16) now proves (F.20.3) for
every prime $\ell$ and for the actual
source collision labels, without requiring
a lift of an arbitrary finite candidate.
The exact maps and finite tests enlarge the
verified winding comparison of F.19.
The remaining full coset equation in
$D^+/K^+$, and the original residual
$A_5$-centralizer, retain their stated scope.
\end{proof}
